\section{Related Work}
\label{sec:related_work}

\subsection{Agent Tool Security and Over-Privilege}
Autonomous agents extend LLM capabilities via external tool use. However, tool-use introduces expanded attack surfaces. Yang et al.~\cite{yang2026when} formalized over-privileged tool selection as a fundamental failure mode across 544 tool scenarios in 8 domains, showing that frontier models frequently ignore lower-privilege tools. Concurrently, \textsc{MiniScope}~\cite{zhu2025miniscope} explored authorization proxies for tool-calling agents to enforce least-privilege without user friction. However, all prior frameworks assume English-centric tool descriptions and user queries, neglecting cross-lingual transfer.

\subsection{Multilingual Robustness in Foundation Models}
Multilingual LLM benchmarks such as \textsc{Maps}~\cite{maps2025} have demonstrated that non-English queries often result in performance degradation and security vulnerabilities in tool-calling agents. In the Southeast Asian context, the \textsc{SEA-LION} initiative~\cite{aisingapore2024sealion} developed culturally and linguistically aligned open-weights models for regional languages. While linguistic benchmarks measure translation and toxicity, security evaluations of agentic tool authorization in regional languages like Javanese remain unaddressed.

\subsection{Runtime Verification and Safe Agent Architecture}
Runtime assurance for autonomous software systems has long relied on formal methods and static type checking. In agentic software repair, \textsc{AutoCodeRover}~\cite{zhang2024autocoderover} employs test-guided validation loops to constrain LLM code generation. Industry security standards such as the OWASP Top 10 for LLM Applications~\cite{owasp2025llm} explicitly highlight the Principle of Least Privilege and deterministic runtime verification as core prerequisites for deploying autonomous agent systems in production.
